Osservatori Standard e le CMF

Poiché ogni persona vede i colori in modo leggermente diverso a causa della sensibilità della propria retina, la scienza ha creato un Osservatore Standard, ovvero un "occhio matematico" medio.

Le CMF (Color-Matching Functions) sono tre curve, chiamate $\overline{x}(\lambda)$, $\overline{y}(\lambda)$ e $\overline{z}(\lambda)$, che descrivono la sensibilità spettrale di questo osservatore ideale.

- Queste curve rappresentano come l'occhio risponde alle componenti del rosso, del verde e del blu per ogni lunghezza d'onda.
- Vengono utilizzate per calcolare i valori di tristimolo X, Y, Z, che sono le coordinate numeriche finali di un colore.

### CIE 1931 e CIE 1964: Qual è la differenza?

La differenza principale tra questi due standard riguarda quanto spazio occupa l'oggetto che stiamo guardando nel nostro campo visivo, espresso come angolo di osservazione.

- CIE 1931 (Osservatore a 2°): È stato progettato studiando la visione umana su un angolo di soli 2 gradi. Immagina di guardare la punta di un dito a circa 50 cm di distanza: in quell'area piccolissima della retina (la fovea) siamo molto sensibili ai dettagli, ma vediamo meno superficie. È adatto per oggetti visti da lontano o molto piccoli.
- CIE 1964 (Osservatore supplementare a 10°): Poiché l'occhio si comporta diversamente quando guarda superfici più ampie, nel 1964 è stato introdotto un osservatore basato su un angolo di 10 gradi. Immagina di guardare un oggetto grande quanto un'arancia a 50 cm di distanza. Questo standard è più accurato per misurare il colore di superfici estese.